The experiments measure, for each of the three tests, the size and the power of the published rule, of the correction and of the variants that isolate each component of the correction. The design lists the predictions first, then the data, the rules compared, the metrics, the repetitions and the statistics. \Cref{tab:design} summarises the four experiments.

\begin{table}[htbp]
\centering
\small
\caption{The four experiments: the engine function each one calls, the seeded generator and the grid of settings it runs, the number of cells and the replications of each cell; every rate of the results carries the standard error of its replications.}
\label{tab:design}
\fitwidth{\begin{tabular}{>{\raggedright\arraybackslash}p{2.4cm}>{\raggedright\arraybackslash}p{3.2cm}>{\raggedright\arraybackslash}p{5.6cm}r>{\raggedright\arraybackslash}p{2.8cm}}
\toprule
Experiment & Engine function & Data and grid & Cells & Replications per cell \\
\midrule
Coverage test & \texttt{conformal\_}\allowbreak\texttt{coverage} & \relax uniform scores, calibration $n$ in 200, 1,000, 5,000; $m/n$ from 0.1 to 8 & 35 & 4,000 \\
Resilience test & \texttt{shift\_}\allowbreak\texttt{development\_}\allowbreak\texttt{reference} & \texttt{CovariateShiftPlant}, evaluation 125 to 1,000, shift 0 to 3 & 10 & 400 \\
Situation light & \texttt{situation\_}\allowbreak\texttt{test} & \texttt{ThresholdBiasPlant}, 400 and 1,000 applicants, gap 0 to 0.8, $k=30$ & 10 & 1,000 (gap 0), 400 (others) \\
Situation light, second design & \texttt{situation\_}\allowbreak\texttt{test} & \relax two independent uniform attributes, 400 and 1,000 applicants, bias 0 to 0.4, $k=30$ & 8 & 2,000 (bias 0), 400 (others) \\
\bottomrule
\end{tabular}}
\end{table}


\subsection{Predictions}

The predictions are bounds on rates. They restate, as testable bounds, the findings of the small simulations of the product's notebook \emph{Measured size and power of the corrected tests}, which ran \NbCovReps{} replications per setting for the coverage test, \NbShiftReps{} for the resilience test and \NbSitReps{} for the situation light. In that notebook the binomial coverage test rejected an exchangeable window in \NbBinOne{} of the replications at equal windows and in \NbBinFour{} at a test window four times the calibration window, against \NbBBOne{} and \NbBBFour{} for the Beta-Binomial test; a fixed drop of 0.05 or 0.10 fired in \NbFixedFiveZ{} and \NbFixedTenZ{} of the shift-free replications against \NbPermZ{} for the permutation reference; and the situation cut 0.10 fired in \NbSitCutTenZ{} of the null replications against \NbSitCutTwentyZ{} for the cut 0.20. Those figures carry large Monte Carlo errors, which the notebook states, and they say nothing about other sample sizes. The experiments here are larger and cover a grid, and the predictions are therefore written as bounds that a larger design can contradict. The predictions were not registered before the experiments: they were written together with the larger experiments, which the version history of the repository shows, since they and the results first appear in the same commit. H1a, H1b, H2a, H2b and H3a restate the notebook's figures as bounds; H2c and H3b to H3e were formulated with the experiments, and H3c, H3d and H3e in particular could have been phrased after the results were seen, so they do not carry the weight of a pre-registered test. The bound used for a size is the nominal level 0.05 plus two Monte Carlo standard errors of the measured rate. \Cref{tab:hypotheses} lists the \HypTotal{} predictions, each with its statement, the value measured and the outcome; \cref{sec:results} reads each in turn. A prediction that the measurements contradict stays in the table as refuted, and none has been removed.

\begin{table}[htbp]
\centering
\small
\caption{The predictions stated before the experiments, each with the value the experiments measured and whether it held: held, held in part, or refuted; the bound used is 0.05 plus two Monte Carlo standard errors of the rate.}
\label{tab:hypotheses}
\fitwidth{\begin{tabular}{>{\raggedright\arraybackslash}p{0.8cm}>{\raggedright\arraybackslash}p{6.6cm}>{\raggedright\arraybackslash}p{4.2cm}>{\raggedright\arraybackslash}p{1.8cm}}
\toprule
Id & Statement & Measured & Outcome \\
\midrule
H1a & \relax The binomial coverage test rejects in at least 0.15 of exchangeable replications at $m/n=1$ and 0.35 at $m/n=4$, up to two standard errors & 0.148 to 0.162 and 0.361 to 0.378 & held \\
H1b & The Beta-Binomial test rejects in at most 0.05 plus two standard errors at every ratio of the grid, and its exact size is at most 0.05 & largest simulated 0.054, largest exact 0.0499 & held \\
H2a & A fixed drop of 0.05 AUC points fires above 0.05 plus two standard errors on shift-free replications; a fixed drop of 0.10 does so at the smallest design & cut 0.05: 0.125 to 0.254 over four sizes; cut 0.10 at 125 rows: 0.098 & held \\
H2b & The permutation reference with 199 draws rejects in at most 0.05 plus two standard errors without shift, at every evaluation size & 0.048 to 0.062 & held \\
H2c & The permutation reference detects a planted shift less often than the fixed cut 0.05 in every cell, and the difference narrows as the evaluation sample grows & shift 2: gap 0.392 at 125 rows, 0.123 at 1,000 rows & held \\
H3a & With 400 applicants and thirty neighbours the amber cut 0.20 fires in at most 0.05 plus two standard errors of the replications without bias and the cut 0.10 above it & cut 0.20: 0.014; cut 0.10: 0.114 & held \\
H3b & The cut 0.20 detects the planted gap of 0.8 in at least 0.95 of the replications & 1.000 & held \\
H3c & The cut 0.20 costs no power against the cut 0.10 at the planted gaps 0.2 and 0.4 & gap 0.4: 0.968 against 1.000; gap 0.2: 0.410 against 0.743 & refuted \\
H3d & The size of the cut 0.10 measured at 400 applicants carries over to 1,000 applicants & 0.114 at 400, 0.018 at 1,000 & refuted \\
H3e & The size of the cuts 0.10 and 0.20 at 400 applicants is the same on a second planted design with independent attributes & cut 0.10: 0.114 and 0.195; cut 0.20: 0.014 and 0.035 & refuted \\
\bottomrule
\end{tabular}}
\end{table}


The first group concerns the coverage test (H1a and H1b): the binomial test rejects far above the level when the test window is not small against the calibration window, and the Beta-Binomial test does not exceed it anywhere on the grid. The second group concerns the resilience test (H2a to H2c): a fixed drop rejects above the level without shift, the permutation reference does not, and the reference detects a planted shift less often than the fixed cut 0.05, by a margin that narrows as the evaluation sample grows. The third group concerns the situation light (H3a to H3e): the cut 0.20 holds the level and the cut 0.10 does not at 400 applicants, the cut 0.20 detects a large planted gap, and three predictions are stated about what the cut costs and whether it transfers: that the cut 0.20 loses no power against the cut 0.10 at small gaps (H3c), that the size of the cut 0.10 measured at 400 applicants holds at 1,000 (H3d), and that the sizes are the same on a second data-generating process (H3e). These three are the predictions on which the correction is weakest.

\subsection{Data and provenance}

Every experiment draws its data from a generator with a known truth, so that no label is observed and every alarm is a false alarm or a detection. The coverage experiment draws calibration and test scores from the uniform law on $(0,1)$; the choice is without loss of generality for continuous scores, because \cref{prop:law} depends on the ranks only. The drift is planted by raising the test score to the power $1/b$, which makes the coverage of a quantile $q$ equal to $q^b$; $b=1$ is the null, and $b$ from 1.02 to 1.20 lowers the mean coverage at the equal windows from the target 0.90 to \CovMeanBTwo{} and \CovMeanBTwenty{}. The calibration sizes are 200, 1,000 and 5,000, and the test window is 0.1, 0.25, 0.5, 1, 2, 4 and 8 times the calibration window; three window pairs, $(20000,2000)$, $(2000,2000)$ and $(1000,4000)$, carry the drift, and they are the pairs of the notebook, so that its figures can be compared with the larger experiment.

The resilience experiment uses the engine plant \texttt{CovariateShiftPlant}: a development sample twice as large as the evaluation sample, a linear model, and a quarter of the evaluation rows (placed at random) shifted by $\delta_s$ standard deviations on the first covariate, in a region where the true logit changes slope, so that the shifted rows are the ones the model ranks badly. The evaluation sizes are 125, 250, 500 and 1,000 rows; the shifts are 0 and 2 at every size, and 1 and 3 at 500 rows. The worst subsample is the share $s=0.25$ of the evaluation rows with the highest resemblance score. The situation experiment uses the engine plant \texttt{ThresholdBiasPlant}: the model decides with a linear predictor, the protected group (half of the applicants) is approved below a cut-off lower by $\gamma$ log-odds, and the reference group is approved at the rate 0.7; the applicants number 400 and 1,000, the gaps are 0, 0.1, 0.2, 0.4 and 0.8, and the test uses $k_{\mathrm{nn}}=30$ neighbours. The second situation experiment uses the design of the engine's own size-and-power suite: two attributes independent and uniform on $(0,10)$, a protected group of share 0.4, and a probability of a negative decision of 0.25 for the others and $0.25+\beta$ for the protected group, independent of the attributes; $\beta$ is 0, 0.1, 0.2 and 0.4.

The committed evidence comes from two documents of the product. The coverage rows are those of the shift variant of the corporate case: a gradient-boosting model fitted on the Polish company bankruptcy data \citep{zieba2016}, a locked holdout of \CmBaseM{} rows, a conformal set calibrated on \CmNcal{} rows, and windows of the holdout shifted by sizes taken from the Federal Reserve hypothetical corporate portfolios. The fairness rows are those of six year pairs of the mortgage decisions of Delaware in the public mortgage data \citep{hmda2018}. The truth of these documents is not planted, so they are used to show what each corrected rule does on the product's data, and never to measure a rate.

\subsection{The rules compared}

All the rules are run by the product on the same replications, through the engine functions or through p-value tables that are checked against them. For the coverage test they are: the binomial p-value (the practice); the Beta-Binomial p-value with the minimum-likelihood rule (the proposed test); the Beta-Binomial p-value with the equal-tail rule, which doubles the smaller tail; a normal test with the variance ratio of \cref{prop:variance}; the Beta-Binomial test with a wrong calibration size, half and double the true one, which measures how sensitive the test is to a mistake in $n$; and the binomial test with its cut calibrated to the exact size, the rule the binomial test would have if its size were corrected by moving the cut. For the resilience test they are: a fixed cut of 0.05 and of 0.10 on the raw drop; a z rule that standardises the difference of the AUC of the rows outside the worst subsample and of the worst subsample with the variance of the structural components of the area \citep{delong1988}; the permutation reference read with the first 19, 39, 99 and 199 of its draws; and a fixed cut at the 95th percentile of the null drops of the same cell. For the situation light they are: the share cut at seven values from 0.05 to 0.30, a one-sided binomial test of the count of significant complainants against the level (which treats the complainants as independent), and the rule that rejects when at least one complainant is significant at the Bonferroni level $\theta/Q$.

\subsection{Metrics, repetitions and statistics}

The metric of every rule is its rejection rate, $\hat r$ of \cref{eq:rate}, with the standard error $\widehat{\mathrm{se}}$ of its replications. The size is the rate under the null; the power is the rate under a planted defect. For the coverage test the exact size of every rule, which has no sampling error, is computed by summing the Beta-Binomial law over the rejection region (\cref{eq:exact-size}), and the exact power by integrating the binomial mass over the Beta law of the coverage; the simulated rates are compared with them in units of the standard error. For the situation light the design effect $\mathrm{DE}$ of \cref{eq:design-effect} describes the null law of the share.

The repetitions are 4,000 per cell for the coverage test (\CovTotalCells{} cells and \CovTotalReps{} replications), \ShReps{} per cell for the resilience test (\ShCells{} cells and \ShTotalReps{} replications) and \SiRepsNull{} at the null and \SiRepsGap{} at each planted gap for the situation light, with \SiInRepsNull{} at the null of the second design. The standard error of a rate near 0.05 is therefore about \CovSeBBB{} for the coverage test, \ShSePermDZ{} for the resilience test and \SiSeCutFifteenSmallZ{} for the situation light at the null. Every replication has its own random stream, built from $\mathrm{SeedSequence}$ with the root seed 20261104, the experiment, the cell and the replication index, so that a replication is reproducible alone and the cells are independent. The experiments are deterministic: a second run from an empty cache reproduces the committed compact results byte for byte.

A prediction is judged by its bound: held when the measured rate is on the stated side of the bound within two standard errors, refuted when it is not, and held in part when some of its clauses hold. A prediction of a difference or a cost (H3c, H3d, H3e) is refuted when the difference exceeds twice the standard error of the difference of two independent rates; H2c holds when the permutation reference has the lower power in every cell with a shift and the gap at shift 2 falls at each larger evaluation size. A comparison of two rules is read on the same replications, and a difference is called real when it exceeds twice the standard error of the difference of two independent rates, which is conservative for paired rules. No correction for multiple comparisons is applied, because the claims are about the ranges of rates over the grids and not about a single cell, and the text names the cells where a difference is within the error.
